% generated by experiments/make_thesis_figures.py -- do not edit by hand
\begin{table}[htbp]
\centering
\footnotesize
\setlength{\tabcolsep}{4pt}
\caption{Noise test: level-matched rule plus the LLM's own residuals (exploratory, post-hoc).}
\label{tab:noise-test}
\begin{tabular}{lcrrcrrc}
\toprule
Noise variant & Post/CPP cloud & LLM & Pct & DH\textsubscript{CPP} cloud & LLM & Pct & Dom. \\
\midrule
\multicolumn{8}{l}{\emph{gemini-3.5-flash-lite (final run v1)}} \\
i.i.d. & 1.238 [1.147, 1.340] & 1.054 & 0 & 155.939 [146.398, 168.636] & 143.8 & 2 & 0/50 \\
i.i.d., orientation only & 1.367 [1.269, 1.474] & 1.054 & 0 & 174.454 [164.503, 181.948] & 143.8 & 0 & 0/50 \\
household-shuffled & 1.110 [1.025, 1.211] & 1.054 & 14 & 148.459 [141.097, 159.793] & 143.8 & 18 & 1/50 \\
\midrule
\multicolumn{8}{l}{\emph{gemini-3.8-flash (Arm A)}} \\
i.i.d. & 1.459 [1.392, 1.525] & 1.261 & 0 & 242.308 [233.210, 251.979] & 240.5 & 34 & 0/50 \\
i.i.d., orientation only & 1.500 [1.422, 1.549] & 1.261 & 0 & 260.841 [251.886, 271.688] & 240.5 & 0 & 0/50 \\
household-shuffled & 1.277 [1.189, 1.408] & 1.261 & 38 & 244.782 [236.580, 256.826] & 240.5 & 14 & 1/50 \\
\bottomrule
\end{tabular}
\par\smallskip
\parbox{\linewidth}{\footnotesize Cloud = median [5th--95th percentile] over 50 seeds. Pct = percentile of the LLM value within the cloud (0 = better than every seed; lower is better on both metrics). Dom.\ = seeds at least as good as the LLM on both post/CPP ratio and CPP-day degree-hours.}
\end{table}
